Validation of user intentions in process orchestration
and choreography
Gerd Gröner a,n
, Mohsen Asadi b
, Bardia Mohabbati b
, Dragan Gašević
c
,
Marko Bošković
d
, Fernando Silva Parreiras e
a WeST Institute, University of Koblenz-Landau, Germany
b
Simon Fraser University, Canada
c
Athabasca University, Canada
d
Research Studios Austria, Austria
e
FUMEC University, Brazil
article info
Available online 6 June 2013
Keywords:
Requirement modeling
Goal-oriented process engineering
Inconsistency detection
Goal-oriented process design
abstract
Goal models and business process models are complementary artifacts for capturing the
requirements and their execution flow in software engineering. In this case, goal models
serve as input for designing business process models. This requires mappings between
both types of models in order to describe which user goals are implemented by which
activities in a business process. Due to the large number of possible relationships among
goals in the goal model and possible control flows of activities, developers struggle with
the challenge of maintaining consistent configurations of both models and their
mappings. Managing these mappings manually is error-prone. In our work, we propose
an automated solution that relies on Description Logics and automated reasoners for
validating mappings that describe the realization of goals by activities in business process
models. The results are the identification of two inconsistency patterns – orchestration
inconsistency and choreography inconsistency – and the development of the corresponding algorithms for detecting these inconsistencies.
& 2013 Elsevier Ltd. All rights reserved.
1. Introduction
With the growing importance of process-aware information
systems (PAISs), business process modeling has gained a
significant research attention [1]. A business process model is
an operational representation of activities, their ordering and
routing in order to achieve goals; that is, delivering services to
customers. However, as outlined in [2], the reasons and
rationales how the execution of a process and the corresponding tasks satisfy the requirements of PAISs are often not
captured within business process models.
Requirements engineering offers proven means for
understanding user intentions. Goal-oriented modeling is
a prominent formalism to describe requirements of a
system in terms of goals (intentions) and relationships
between goals. A goal describes a certain system functionality or property that should be achieved. Thus, it is not
surprising that the recent research on PAISs has focused on
the integration of business process modeling with goaloriented modeling (cf. Section 8). Related research concentrates on basic mapping and aligning principles between
the intentional perspective, represented by goal models
[3–5] or business models [6,7] and the process model
perspective. The main focus of these approaches is either
on enriching a process model with goals and relationships
between goals or on transformations from goal models to
process models. However, less attention has been paid
to the analysis and formalization of the influence of
intentional relations, which reflect the user requirement
Contents lists available at ScienceDirect
journal homepage: www.elsevier.com/locate/infosys
Information Systems
0306-4379/$ - see front matter & 2013 Elsevier Ltd. All rights reserved.
http://dx.doi.org/10.1016/j.is.2013.05.006
n
Corresponding author. Tel.: +49 261 287 2447.
E-mail addresses: groener@uni-koblenz.de (G. Gröner),
masadi@sfu.ca (M. Asadi), mohabbati@sfu.ca (B. Mohabbati),
dragang@athabascau.ca (D. Gašević),
boskovic@researchstudios.at (M. Bošković),
fernando.parreiras@fumec.br (F. Silva Parreiras).
Information Systems 43 (2014) 83–99
perspective, on process model relations, which cover the
process control flow and message exchange between
processes.
In this paper, we tackle the challenge of formalizing the
realization of intentions by activities in business processes.
These realizations are described by mappings. Based on the
And
Process
Order
Minimize Risk
Customer
Satisfaction
Bill, Build,
Then Ship
+
And
Build & Package
Order
Electronic
Payment
In Person
Payment
CP
EP IPP
or
Ship & Bill
Apply
Process to
Customer
Or
+
Determine
Trustworthiness of
Customer
And
Check
Credit Rate
Check if Return
Customer
DTC
CCR CRC
Xor
Approve Order
AO
Apply Process
to Any
Customer
And
BPO Apply Process to
Trusted Customer
Ship Order Bill Preparation
And
Payment
Collect
Payment
Apply
Discount
+
Customer Profile
Retrieve
CPR
SO BL
AD
Ship
ment
Front
Store
Back
Store
Supplier
Continuing
Supply
Buy Item
Items
And
Ship Order
Continuing
Business
Place Order
Search
Keyword
Browse
Catalogue
Process Orders
Custo
mer
Handel
Payment
Goal
Softgoal
Task
Decomposition
Contribution
Dependency
+
+
+ +?
Hurt Some
Negative
Make Some
Positive
Help Unknown
Break
Actor
[ BL]
Store Customer
Send
Credential
Information
Custom er
Identification
Submit the
Order
Ask for order
Processing
Credit Card
Checking
Custom er
Record
Checking
Custom er
Profile Retrieve
Approve Order
Ask for
Paym ent
Method
Select
Paym ent
Method
Fraud
Detection
Credit Card
Debit Card
Money Order
Verify Paym ent
Information
Apply Discount
Build and
Package Order
Shipment
Send Receipt
Receive I tem s
Receive
Receipt
[ AO]
[ CCR]
[ CRC]
[ CPR] [ I PP]
[ EP V
I PP]
[ AD]
[ BPO]
[ SO]
[ CP] [ CP]
Choreography I nconsistency Strong Inconsistency
Potential I nconsistency
Gateway
AND OR XOR
Mapping
[ …]
(goal to process activity)
Legend
Fig. 1. Goal realizations in business process models. (a) Goal model of the E-Store. (b) Business process model of the E-Store.
84 G. Gröner et al. / Information Systems 43 (2014) 83–99
formalization, we specify orchestration and choreography
inconsistencies that might occur due to contradicting relationships between activities and their corresponding
(mapped) intentions. In that sense, an automated validation
of the mappings needs to assure the fulfillment of user goals
in each execution configuration/path of business processes.
To address this challenge, our first contribution (Section
4) is the definition of types of inconsistencies between goal
and process models. The definition of inconsistencies is
based on correspondences between intentional relationships of the goal model and relationships in the process
model, which are given by control flow relations (orchestration) in terms of workflow patterns [8] and by interaction
relations (choreography) between processes. Next, based on
these correspondences, we propose a formal modeling
(Section 5) and validation (Section 6) approach in Description Logics. The focus of the modeling is to capture the
relationships among activities of the business process
model and between user intentions of the goal model.
Sound and complete Description Logics based reasoning is
used to ensure realization equivalence between user intentions and their corresponding (mapped) activities; that is,
that there are no contradictions between relations of
intentions compared to their mapped activities. Finally, we
compare our approach to already available approaches for
validation of correspondences between goal and process
models (Section 8) and outline the conclusions (Section 9).
2. Foundations
This section presents fundamentals of the requirement
perspective and the business process design perspective.
2.1. Goal models
To date, several languages for goal modeling have been
proposed, e.g., i*/Tropos [9,10], NFR [11], KAOS [12] and Goal
Requirements Language (GRL) [13]. We adopted GRL, a part
of the recent Recommendation of the International Telecommunications Union named User Requirements Notation
(URN) [13]. GRL is a language that integrates core concepts
of i* and NFR [14].
A goal model consists of actors, intentional elements, and
links [13].
Intentional elements are (hard) goals, soft goals, tasks, and
resources. A (hard) goal is a state of affairs that an actor
would like to achieve. Soft goals are similar to hard goals,
but without a clear-cut criteria if the condition is achieved
or not. Tasks specify conceptual solutions in order to achieve
a particular intentional element. Resources are informational or physical entities, which are available to the actors.
Actors are entities and refer to stakeholders or systems.
Actors can have intentions. In this case, we say an actor has
a scope and activities of this actor are within the
actor's scope.
A concrete goal model is depicted in Fig. 1(a). The goal
model is part of an online shopping case-study, which is
used throughout the paper. The goal model contains five
actors, namely Front Store, Back Store, Customer, Shipment,
and Supplier. When representing user intentions for a
system or a part of a system, actors are either internal or
external actors. Internal actors are part of the system-to-be,
while external actors represent stakeholders that interact
with system actors. In the goal model in Fig. 1(a), only the
intentional elements within the scope of the internal actor
Back Store are depicted. Intentional elements in the scope
of readability.
We distinguish between two kinds of links. (1) Dependency links connect intentional elements and/or actors. (2)
Decompositions and contributions link intentional elements.1
Dependencies describe inter-actor relationships. They
express how a source actor (depender) depends on a
destination actor (dependee) for an intentional element
(dependum) [13]. Dependency relations are only defined
between different actors and/or intentional elements of
different actors, i.e., intentional elements in the scope of
different actors. Hence, a dependency link is not defined
between intentional elements of one actor. According to the
type of the dependum, dependencies are called task, (hard)
goal, soft goal and resource dependencies [15]. For example,
Front Store actor depends on Back Store actor for achieving
Process Order objective, and the Supplier actor depends on
Back Store actor for the soft goal Continuing Business.
Decompositions specify which intentional elements a
certain intentional element needs for its satisfaction. For
instance, the goal Ship & Bill is satisfied if both tasks Ship
Order and Bill Preparation are fulfilled. GRL supports AND,
IOR and XOR decompositions. An AND decomposition specifies that all source intentional elements need to be
satisfied for the target intentional element to be satisfied.
IOR is used to specify that the satisfaction of at least one
source satisfies the target, whereby XOR specifies that
exactly one of the source elements is necessary to satisfy
the target.
Contributions express how an intention contributes to
the achievement/fulfillment of another intention. Fig. 1(a)
shows different contribution types. The Make and Break
contributions are respectively positive and negative, and
sufficient for the fulfillment of a target element. Help and
Hurt contributions are also positive and negative respectively, but insufficient. The extent of the contribution of
SomePositive and SomeNegative is unknown. Finally, for
the Unknown contribution link, both the extent and degree
(positive or negative) of the contribution are unknown.
Definition 2 compactly states the specification of goal
models.
Definition 1 (Goal model). A goal model is a quintuple
GM ¼ðA; G;P; D; CÞ. A denotes a set of actors. G is a set of
intentional elements (also called intentions or goals). Intentions are (hard) goals (Gg), tasks (Gt), soft goals (Gs), and
resources(Gr).
 P is a set of dependency links:
PDðA∪GÞ G ðA∪GÞ:
1 According to the GRL standard [13] it is also possible to specify
correlation links. Correlation links are similar to contribution links, except
that they model side effects. From the goal satisfaction perspective both of
these links are same. For this reason, we restrict ourselves only to
contribution links.
G. Gröner et al. / Information Systems 43 (2014) 83–99 85
 D are decompositions of intentional elements:
DDG fAND; IOR; XORg PðGÞ:
2
 C denotes positive and negative contributions:
CDG  fþ

; þ

; þ; 7?; 

;−; þ

g G:
Regarding this definition, dependencies ðPDðA∪GÞ
G ðA∪GÞÞ are relations between intentional elements ðGÞ
within the scope of different actors or between actors ðAÞ.
For instance in Fig. 1(a), there is a dependency from actor
Front Store to hard goal Process Order, which is in the scope
of actor Back Store. The second argument ðGÞ refers to the
dependum, e.g., the hard goal Process Orders in Fig. 1(a).
In essence, goal models represent user intentions, relationships between intentions, actors, which might have
intentions within their scopes, and dependencies between
actors.
2.2. Business process models
There is an emergence and proliferation of processoriented software development methods for organizations
and enterprises. Today, software products as services are
designed, built, executed, maintained and evolved by means
of business processes on top of Web services technologies
[1,16]. To accomplish this, business process modeling distinguishes between the concepts of orchestration and choreography [17].
A process orchestration describes the control flow of
activities in the perspective of one party. It refers to an
executable part of a business process that is provided by one
party through a central coordination. Activities realize and
implement certain intentions that can be specified by a
requirement model [18].
Process choreography takes inter- and intraorganizational processes and communication between
them (e.g., B2B and B2C) into account. This may involve
multiple parties that perform different parts of the overall
business process. Choreography describes the set and ordering of interactions (i.e., message exchanges) between individual parties and defines the expected behavior among the
interacting parties, i.e., a type of procedural contract or
protocol.
2.2.1. Process orchestration
Workflow patterns [19] describe the orchestration of
processes. They define how the process flow proceeds in
sequences and splits into branches and how branches
converge.
Business process models can be designed at different
levels of abstraction. A number of graph-based business
process modeling languages have been proposed, e.g.,
Event-driven Process Chain (EPC), Yet Another Workflow
Language (YAWL) [20] and UML Activity Diagrams. We use
the Business Process Modeling Notation (BPMN) [21] as a
base and standard modeling language. Nonetheless, the
proposed approach given in the present work is independent of any goal and process modeling languages. Despite a
variance of modeling notations and different semantics, all
modeling languages share the common concepts of activities, gateways (or routing nodes) and relationships
between them.
Definition 2 (Business process model). A business process
model is defined as a connected graph ΓPM ¼ðV; EÞ, where V
depicts a set of vertices, denoting activities A and gateways
G (V ¼ A⊎G). A gateway G∈G has a type TðGÞ such that
TðGÞ∈fAND; IOR; XOR;DISCg. EDV  V denotes a set of edges
between vertices.
To allow the analysis of activities and their relationships
with intentions and their corresponding intentional relationships, we focus on structured process models, which are
more comprehensible and less error-prone [22]. We consider structural well-formedness conditions for business
process models [23,24], where each process model is
decomposed into aggregated SESE (single-entry-single-exit)
fragments. SESE fragments can be derived from a process
model in linear time (cf. [25]). According to these structural
well-formedness conditions, gateways can have either
exactly one incoming edge and multiple outgoing edges or
multiple incoming edges and exactly one outgoing edge.
Furthermore, we assume a materialized representation
according to Definition 3.
Definition 3 (Materialized business process model). Given a
business process model ΓPM ¼ðV; EÞ, a materialized business
process model PM is a quintuple ðA; G; E; F; EAÞ. Vertices V
are either activities A or gateways GðV ¼ A⊎GÞ. The set
F DV  V  B represents single-entry-single-exit (SESE)
fragments, in which B represents a set of branches between
entry and exit vertex. A branch B∈B is considered as a set of
atomic activities. EADE
þ
is a set of sequence edges obtained
from the process graph by treating gateways as
transparent.3
The set EA captures the explicit predecessor and successor relationships between atomic activities without gateways. This set is derived from the set of edges E by
neglecting gateways and replacing them by their corresponding next predecessor or successor activity. A business
process model is illustrated in Fig. 1(b). Mappings between
activities and tasks in the goal model are indicated by
activity annotations. E.g., the activity Credit Card Checking
is mapped to goal Check Credit Rate (CCR).
Fragments F DV  V  B (as introduced in Definition 3)
can be found for instance in the Customer process on the
right side, where AND gateways are the start and end point
of the fragment and B consists of two branches (between
the opening and closing gateway). Each branch B∈B contains only one activity (Receive Items and Receive Receipt).
2 PðGÞ denotes the power set of G.
3
The set E
þ denotes the transitive closure on the set of edges E. EA
denotes the subset of the transitive closure such that activities are either
directly connected or transitively connected, but only with gateways inbetween them.
86 G. Gröner et al. / Information Systems 43 (2014) 83–99
In BPMN, process orchestrations are represented within
container called pools. In Fig. 1(b), there are two pools:
Store and Customer. Both pools contain a business process.
A pool refers to a participant and its local process view. In
the pool Store, the process is spread over two lanes
(FrontStore and BackStore). A lane is a sub-partition of a
process.
2.2.2. Process choreography
A choreography defines the sequence of interactions
between participants in terms of message exchange, which
appears between pools in BPMN models. A choreography
does not exist within a single pool (i.e., in a local context of
control). Fig. 1(b) illustrates an example of a choreography
between the pools Store and Customer. For instance, there
is a message sent from the activity Submit the Order in the
pool Customer to the activity Ask for Order Processing in the
pool Store. Message flows between pools define a choreography [26]. We consider a process choreography as a set of
materialized business process models and a set of message
flows between activities, formally defined as follow.
Definition 4 (Process choreography). Given a set P of
materialized business process models. The set MD ^A  ^A,
with ^A ¼ ⋃P∈P∧A ¼ activitiesðPÞA is a set of message flows. A
process choreography is a tuple Ch ¼ðP;MÞ with the condition:
∀ðA1; A2Þ∈M : ⇒A1∈activitiesðP1Þ∧A2∈activitiesðP2Þ∧P1≠P2.
4
For each message exchange between A1 and
A2ððA1; A2Þ∈MÞ it is required that A1 and A2 are activities of
different processes. We require that each process is in a pool
and each pool contains one process.
3. Methodology overview
This section provides an overview of the proposed
methodology and generated artifacts. This methodology
consists of two subprocesses: (i) the development of process
models based on intentions and (ii) the validation, as illustrated in Fig. 2.
In the former subprocess (see Fig. 2(a)), requirement
engineers develop a goal model representing stakeholders'
intentions and their refinement into tasks. This procedure
starts with identifying the actors (e.g., Back Store and Front
Store) of the system-to-be and their high level goals (e.g.,
Process Order) and soft goals (e.g., Minimize Risk). Actors
dependencies are identified and modeled in the goal model.
High level goals of actors are decomposed by following the
framework proposed by Liaskos et al. [27] where intentional
variability concerns are recognized for each high level goal.
Then, these goals are refined according to the variability
concerns. After refining goals, requirements engineers analyze the impacts of goals on other goals and model these
impacts using contribution links. Afterwards, a process
model is produced by defining activities and subprocesses
that realize tasks of the goal model. Activities are organized
in swim-lanes. Each lane refers to an internal actor in the
goal model. According to the relations in the goal model,
process orchestration and choreography relations are determined and represented. Simultaneously, mappings between
tasks in the goal model and activities in the process model
are devised. These models offer two different viewpoints
(requirements on one side and process orchestration and
choreography on the other side), covering artifacts and
relationships among artifacts for different purposes.
The second subprocess (see Fig. 2(b)) identifies possible
inconsistencies that may happen between relationships in
the goal model and process models. As a first step, the goal
model, the process model and the mapping model are
transformed into a knowledge base in Description Logics.
The knowledge base covers the relationships of elements
from both models. Afterwards, the second step leverages
reasoning services, which rely on the well defined Tarskistyle semantics [28] of Description Logics, to recognize
(possible) inconsistencies between goal models and process
models. To this end, a subset of Description Logics constructs (ALC expressiveness) is used in our framework.
4. Realization inconsistencies
Tasks in the goal model are those intentional elements
that need to be performed in order to achieve the specified
requirements. The basis for aligning business process models with goal models is to map tasks of the goal model to
activities (atomic activities or subprocesses) in the business
process model in order to express the implementation of a
goal by an activity. We cover the following mapping rules
between elements in the goal models and business process
models, which follow basic mapping and realization principles in the literature [3,29,30]:
 Actor mapping: The actors within a goal model are
mapped to parties that execute activities in the business
process, i.e., to swim-lanes and pools. In our case, each
pool contains one process.
○ Internal actors (actors that belong to the system-tobe) are mapped to the lanes of a pool. If a pool
consists of multiple lanes the activities of a process
(in this pool) might be distributed over the lanes of
the pool.
○ External actors (actors that interact with the systemto-be) are mapped to pools.

Intentional element mapping: Among the intentional
elements, tasks are implemented by activities in a
process. This is expressed by mappings.
○ Tasks in the goal model are mapped to the activities
(either atomic or composite, i.e., sub-processes).
These mapping principles follow the modeling intuition
of both models: (i) Actors in a goal model represent active
entities that perform actions to achieve their goals. These
goals are within the actor's scope. Likewise in the business
process model, lanes group activities of a process and pools
group processes that refer to different roles. (ii) Tasks in the
goal model are achieved by executing activities (atomic or
composite activities) in a process.
4
Ai∈activitiesðPi
Þ is an abbreviation to denote that activity Ai∈A occurs
in process Pi ¼ðA; G; E; F; EAÞ.
G. Gröner et al. / Information Systems 43 (2014) 83–99 87
In Fig. 1, mappings are represented by annotations of
activities with the name of intentional elements. For
instance, the activity Money Order is mapped to the intentional element In Person Payment (IPP).
When mappings are established, the process orchestration (ordering of activities) must be consistent with intentional relations between goals, and a process choreography
(relations that are given by message exchange) must exist
between pools and lanes in case there are dependencies
between the corresponding actors.
4.1. Orchestration realization inconsistencies
The process orchestration is described by control flow
relations between activities, called workflow patterns.5
Definition 5 specifies workflow patterns as a collection of
all control flow relations an activity depends on.
Definition 5 (Workflow patterns). Let PM ¼ðA; G; E; F; EAÞ
be a materialized business process model. We denote with
WFA the workflow relations of an activity A ∈A, specified on
activities A1; …; An∈A.
The workflow patterns of a process describe activities and
their execution ordering. The execution ordering is specified
by a combination of the introduced modeling constructs:
activities, gateways and edges (flow edges) between activities and gateways. The relations can be grouped into (i)
basic control flow patterns, (ii) advanced branching and
synchronization patterns, (iii) structural patterns, (iv) multiple instance patterns, (v) state-based patterns and (vi)
cancelation patterns.
Basic control flow patterns cover the elementary aspects
of a process like sequences, parallel and exclusive branching. Advanced branching and synchronization patterns offer
more complex branching and merging constructs like
multi-choice, discriminator and n-out-of-m joins. Structural
patterns contain restrictions regarding the structure like
loops with multiple entry and exit points and implicit
termination. Multiple instance patterns represent processes
where certain activities can have multiple instances (within
one process instance). State-based patterns allow the specification that the process execution is determined by the
state of the process instance at runtime. The cancelation
pattern is used to describe possibilities and situations
where the process execution can be terminated in certain
circumstances.
According to the mapping principles, activities implement tasks of the goal model. These tasks might depend on
other intentional elements in the goal model. This is
described by intentional relations (Definition 6).
Definition 6 (Intentional relations). Let GM ¼ðA; G, P; D; CÞ
be a goal model. IRG denote the intentional relations of the
intentional element G ∈G on other intentional elements
G1;…; Gm∈G. Intentional relations are represented by
decompositions D, contributions C and dependencies P
between intentional elements in the scope of different
actors.
According to the mapping rules, tasks in the scope of each
actor are mapped to activities inside pools and lanes. Therefore, the execution relations between activities (i.e., workflow patterns) in business process models must be consistent
with intentional relations between intentional elements in
goal models. Mappings “carry” intentional relations to the
business process model. Thus, there might be inconsistencies
between the relations given by the workflow pattern of an
activity and the corresponding intentional relations of its
mapped intentional element. We distinguish between two
types of orchestration inconsistencies: (1) strong inconsistency and (2) potential inconsistency (cf. Definitions 7 and 8).
Fig. 2. Goal realizations in business process models. (a) Steps for developing the goal model, process models and mapping model. (b) Steps for detecting
inconsistencies.
5
http://www.workflowpatterns.com.
88 G. Gröner et al. / Information Systems 43 (2014) 83–99
Definition 7 (Strong inconsistency). Assume a workflow
pattern WFA for an activity A ∈A is specified over activities
A1;…; An∈A and an intentional relation IRG for G∈G, defined
over intentional elements G1;…; Gm∈G. Activities
A; A1; …; An are realizations of intentional elements
G; G1; …; Gm. A strong inconsistency between WFA and IRG
occurs if there is no execution combination of activities that
leads to the fulfillment of the intentional relation IRG.
A strong inconsistency says that for no execution that is
expressed by workflow pattern/relation WFA there is a goal
fulfillment/achievement of IRG possible. A weaker notion is
the potential inconsistency. There can exist an allowed
execution WFA in which the corresponding intentional
relation IRG is not fulfilled.
Definition 8 (Potential inconsistency). Assume a workflow
pattern WFA for an activity A∈A is specified over activities
A1;…; An∈A and an intentional relation IRG for an intention
G is defined over intentional elements G1;…; Gm∈G. Activities A; A1; …; An are realizations of intentional elements
G; G1; …; Gm. We define a potential inconsistency between
WFA and IRG if some execution combinations of activities
lead to the fulfillment of the intentional relation IRG and
some execution combinations of activities do not lead to the
fulfillment of IRG.
Both kinds of inconsistencies are recognized based on a
set of existing mappings. Adding additional mappings might
lead to further inconsistencies.
As our approach is applied for process engineering and
modeling at design time, i.e., we consider process models
rather than executions and execution observations, various
patterns of the multiple instance, state-based and cancelation patterns are not affected by in our approach, since they
deal with rather run-time related aspects that cannot
become inconsistent at design time. Table 1 summarizes
the influence of intentional relations on workflow patterns,
but it also provides a comprehensive overview how both
kinds of relations can be mapped to each other. The last
column in Table 1 depicts the influence of dependencies in
the goal model. In this case, we require that the corresponding intentions are in the scope of different actors.
The example in Fig. 1 contains a strong and a potential
inconsistency. The strong inconsistency is due to the activities Send Receipt and Shipment that are mapped to tasks
Bill (BL) and Ship Order (SO), whereby these tasks are ANDsiblings. Thus, each satisfaction of the target element Ship &
Bill requires that both tasks Bill (BL) and Ship Order (SO) are
fulfilled simultaneously, while each process execution
allows either the execution of Send Bill or Shipment.
A potential inconsistency is caused by the mapping of
activities Credit Card Checking and Customer Record Checking to the exclusive sibling tasks Check Credit Rate (CCR)
and Check if Return Customer (CRC) in the goal model.
There are executions where both activities are executed, but
this would contradict to the exclusiveness of the tasks
Check Credit Rate (CCR) and Check if Return Customer
(CRC) to fulfill their target goal DTC.
As shown in Table 1, we use four symbols to indicate
whether we can identify an inconsistency or not: (i) ‘✓’
indicates no inconsistency between intentional relations IR
and workflow relations WP, (ii) strong inconsistencies are
denoted by ‘ ’, (iii) the symbol ‘7’ refers to potential
inconsistencies, and (iv) if our approach cannot lead to
inconsistencies at the modeling level based on a mapping
between tasks and activities, the corresponding cells are
marked with ‘–’. Inconsistencies are caused by contradicting
relationships of mapped elements. Thus, adding additional
mappings might cause further inconsistencies.
4.2. Choreography realization inconsistencies
Choreography describes the message exchange between
processes. Recurrent business process choreography scenarios are formulated using service interaction patterns [31].
These patterns encompass a variety of interaction scenarios
Table 1
Correspondences and mapping influence between workflow patterns and intentional relations.
Workflow patterns Intentional relations
AND IOR XOR þ
 þ
 Dep.
Sequence ✓✓ ✓ ✓
AND–AND parallel split–synchronization ✓✓ ✓ ✓
AND–IOR parallel split–multi-merge ✓✓ ✓ ✓
AND–DISC parallel split–discriminator ✓✓ ✓ ✓
AND–XOR parallel split–simple merge ✓✓ ✓ ✓
XOR–XOR exclusive–simple merge ✓✓ ✓
IOR–IOR multi-choice–synchronizing merge 7 ✓ 7 777
IOR–XOR multi-choice–simple merge 7 ✓ 7 777
IOR–DISC multi-choice–discriminator 7 ✓ 7 777
Arbitrary cycles – – – –––
Implicit termination – – – –––
Multiple instances (MI) pattern – – – –––
Deferred choice 7 7 ✓
Interleaved parallel routing ✓ 7 7
Milestone – – – –––
Cancelation pattern – – – –––
G. Gröner et al. / Information Systems 43 (2014) 83–99 89
ranging from simple message exchanges to multiple participants and multiple message exchanges.
Based on the number of involved parties and maximum
message exchange between two parties, the interaction
patterns are classified to:
 single-transmission bilateral interaction patterns: send,
receive, send/receive;
 single-transmission multilateral interaction patterns:
racing incoming messages, one-to-many send, onefrom-many receive, on-to-many send/receive;
 multi-transmission interaction patterns: multi-responses,
contingent requests, atomic multicast notification.
According to the terminology and modeling assumptions
of the BPMN language, choreography is between processes
of different pools. Thus, the following subset of the interaction patterns is supported in BPMN [32]: Singletransmission bilateral interaction patterns cover the message exchange between two participants, including (i.a)
send, (i.b) receive and (i.c) send and receive. The single
transmission multilateral interaction pattern describes message exchange between one participant on one side and
multiple participants on the other side. In the realm of the
BPMN language, we consider the pattern (ii.a) racing incoming messages. Finally, multi-transmission interaction pattern
allows multiple rounds of message exchange between two
actors. The BPMN language supports the (iii.a) multiresponses pattern.
When intentional elements are mapped to activities,
there might be inconsistencies between the actor dependencies in the goal model and the message exchange
between pools (i.e., participants in an interaction).
In the goal model, dependency links are utilized to
represent the interactions between actors. According to
Mahfouz et al. [30], actors' dependencies can be classified
based on the types of dependum and the logical and
physical nature of the dependency. The dependum element
can be either a goal, a task, or a resource [15]. A dependency
can be either physical—a physical occurrence that indicates
the satisfaction of a dependency (e.g., shipping items into a
warehouse) or informational—needed information are provided for a depender by a dependee. We follow the
argumentation of [30] that physical activities, which participants perform, are not necessarily manifested in the
choreography description in a direct way since the satisfaction of the dependency can only be assessed if the depender
has observed a physical occurrence of indicators. Therefore,
we focus on the informational dependencies.
Interactions between actors are realized in the corresponding business processes in two ways. Firstly, if actors
are mapped to pools, dependencies are realized through
message exchange between pools, or more precisely
between the processes in the different pools. Secondly, in
the case the actors are mapped to lanes, dependencies
between intentions (of different actors) are realized through
control flow relations [3,29]. Therefore, there must exist a
service interaction pattern between activities that are
mapped to the intentional elements of a depender and
activities that are mapped to the intentional elements of a
dependee. Accordingly, there might be a choreography
inconsistency in the business process model with respect
to actor dependencies. A choreography inconsistency is
defined as the following:
Definition 9 (Choreography inconsistency). Let GM ¼
ðA; G;P; D; CÞ be a goal model and G1∈G, G2∈G be intentional
elements, Act1; Act2∈A actors and G1 is in the scope of actor
Act1ðG1∈ScopeðAct1ÞÞ and goal G2 is in the scope of actor
Act2ðG2∈ScopeðAct2ÞÞ.
Let Ch ¼ðP;MÞ be a process choreography. Assume there
are mappings between intentional elements G1 and
A1∈activitiesðP1Þ and between G2 and A2∈activitiesðP2Þ
(P1; P2∈P). Furthermore, there is a dependency relation
ðG1; Gx; G2Þ∈P (G1 is the depender, G2 is the dependee and
Gx is the dependum).6
There is a choreography inconsistency if there is no
service interaction patterns between activities A1 and A2
(i.e., ðA1; A2Þ∉M and ðA2; A1Þ∉M).
Fig. 1 shows an example of a choreography inconsistency. We map Back Store, Front Store, and Customer to the
corresponding swim-lanes (BackStore, FrontStore) in the
process Store and to the process Customer. The actor
Customer depends on the actor Back Store for achieving
the goal Handle Payment. More precisely, the intentional
element Collect Payment (CP) (scope of actor Back Store)
depends on Select Payment Method (scope of actor Customer).7
The task Collect Payment (CP) (scope of the actor
Back Store) is mapped to the activity Ask for Payment
Method in the Store process (FrontStore lane). As we can
see, there is no interaction (i.e., message flow pattern)
between the corresponding activities in the processes for
realizing this dependency. Therefore, there exists a choreography inconsistency since we expect a message exchange
between the activity Ask for Payment Method and the
corresponding activity Select Payment Method (Customer
process).
5. Knowledge base for realization validation
We use Description Logics (DL) [33] to model workflow
patterns, intentional relations and mappings. Based on this
representation, DL reasoning services are used to validate
realizations of intentional elements by workflow patterns.
5.1. Foundations of Description Logics
DL is a decidable subset of first-order logic (FOL). A DL
knowledge base consists of a TBox (Terminological Box) and
an ABox (Assertional Box). The TBox is used to specify
concepts, which denote sets of individuals and roles defining binary relations between individuals. The main syntactic
constructs are depicted in Table 2, supplemented by the
corresponding FOL expressions. Concept inclusion axioms
6
From a logical point of view, a dependency relation depicts a
relationship between actors (depender and dependee) or between intentions of their scope, but there is no relationship imposed to the dependum.
7
Remember that in Fig. 1(a) the intentional elements of the actor
Customer are omitted, thus, the intentional element Select Payment
Method is not depicted in Fig. 1(a).
90 G. Gröner et al. / Information Systems 43 (2014) 83–99
C⊑D mean that each individual of the concept C is also an
individual of D. There are two special concepts in Description Logics, namely the universal concept (top concept) ⊤
and the bottom concept ⊥. The top concept ⊤ is the superconcept of all concepts, i.e., C⊑⊤ holds for each concept C.
⊥ is an unsatisfiable concept. A concept equivalence (or
definition) C≡D is an abbreviation for two concept inclusion
axioms C⊑D and D⊑C. R and S denote roles (properties),
which are binary relations between concepts. R⊑S describes
that R is a subrole of S.
A concept union is a complex concept expression and
refers to a disjunction in FOL, i.e., an individual of an
concept union has to be an individual of at least one concept
Ci of the union. Likewise, a concept intersection refers to a
conjunction in FOL. A concept negation :C is the set of all
individuals that are not individuals of the concept C. C can
be an arbitrary complex concept.
Universal and existential quantifiers either restrict or
require relations with other concepts (C in Table 2). These
other referenced concepts can be arbitrary complex
concepts.
Inference services of DL rely on the well defined Tarskistyle semantics [28]. The subset of DL constructs we use in
our models (ALC expressiveness) in combination with
existing highly optimized reasoning algorithms and systems
allows for practically efficient reasoning. In the remainder of
this paper, we use subsumption checking and concept
classification as basic reasoning services.
5.2. Representation of models and realizations
The key part of our modeling formalism contains the
relations of both models, combined with mappings between
them. The goal model describes intentional relations
between goals and the business process model specify
control flow relations on activities, which refer to basic
workflow patterns [8,19]. According to the inconsistency
detection problem, as specified in Section 4, our approach
is based on a comparison of these relations. In the remainder of this section, we consider how such a comparison is
realized by concept comparison (i.e., subsumption checking)
in Description Logics.
5.2.1. Representation of intentional relations
The intentional relations of a goal model GM are
described in a DL knowledge base ΣGM. Algorithm 1 depicts
the representation of intentional relations of a goal model
GM ¼ðA; G;P; D; CÞ. Only tasks are implemented by activities. Thus, the algorithm introduces for each task G (G∈Gt)a
(complex) concept RelG to capture its intentional relations,
which are either decompositions or contributions. Complex
concepts describe relationships of atomic concepts and
roles (properties) in DL. Intentional elements G are represented as atomic concepts.
Algorithm 1. Intentional relations ΣGM.
1: Input: Goal model GM ¼ðA; G;P; D; CÞ
2: for all A∈A do
3: if ðG; IOR;fG1;…; GngÞ∈DA ðDA DDÞ then
4: RelGi≡⊔j ¼ 1;…;n∃requires:Gj
ðfor i ¼ 1; …; nÞ
5: end if
6: if ðG; AND;fG1;…; GngÞ∈DA ðDA DDÞ then
7: RelGi≡⊓j ¼ 1;…;n∃requires:Gj
ðfor i ¼ 1;…; nÞ
8: end if
9: if ðG; XOR;fG1; …;GngÞ∈DA ðDA DDÞ then
10: RelGi≡⊗j ¼ 1;…;n∃requires:Gj
11: end if
12: end for
13: for all ðG′; þ

;GÞ∈C do
14: RelG≔RelG⊓∃requires:G′
15: end for
16: for all ðG′; þ

;GÞ∈Cdo
17: RelG≔RelG⊓:∃requires:G′
18: end for
19: for all ðGdepender; Gdependum; GdependeeÞ∈P do
20: RelGdepender ≔RelGdepender ⊓∃requires:Gdependum
21: DepGdepender ≔DepGdepender ⊓∃requires:Gdependum
22: end for
Lines 3–5 treat IOR-decompositions of an intention into
subgoals Gi
i∈ð1; …; nÞ. Thus, intentions Gi are disjunctively
related to each other. This is represented in DL by a concept
union over all intentions Gj
that are members of the IORdecomposition. For each intention Gi
that is part of the IORdecomposition, we introduce a concept RelGi
to describe the
relations of each intention. In this vein, a conjunctive
decomposition (lines 6–8) is described by a concept intersection in DL. As in the previous case, we introduce relation
concepts RelGi
to capture conjunction for all members of the
decomposition.
The representation of exclusive decompositions is
straightforward (lines 9–11).8 An intention can only be the
source of one decomposition, i.e., either IOR, XOR or AND.
Sufficient positive contributions (lines 13–15) specify
that the fulfillment of G requires the fulfillment of G′. Thus,
we add the expression ∃requires:G′ to the definition of the
relation concept RelG. The meaning in Description Logics is
that an instance of goal G′ must have a role / property
requires to another goal instance. Sufficient negative contributions (lines 16–18) use concept negation in order to
represent that the fulfillment of G cannot be achieved if G′ is
fulfilled. A task might be involved in multiple (positive and
negative) contributions simultaneously. Contributions that
are not sufficient are not included in this representation
since their effect cannot be completely determined at
design time when mappings are assigned to the models.
Lines 19–22 capture dependencies between intentions,
where intention Gdepender depends on the intention Gdependee
on an intention Gdependum. Logically, this is an implication
that the fulfillment of Gdepender requires the fulfillment of the
intention Gdependum, which is reflected in the DL concept
expression RelGdepender . Following the discussion in Section 4,
this relation is only added into the knowledge base if there
is an explicit dependency relation between the intentional
elements Gdepender and Gdependee and both are in the scope of
different internal actors. In line 21 the dependency between
8 Note that ⊗ is not a standard operator in DL. For a more concise
representation, we use ⊗G′∈fG1 ;…;Gn g∃requires:G′ as an abbreviation for
⊔G′∈fG1 ;…;Gn g∃requires:G′⊓:ð⊔G″
; G‴∈fG1; …;Gngð∃requires:G″⊓∃requires:G‴ÞÞ.
G. Gröner et al. / Information Systems 43 (2014) 83–99 91
intentions Gdepender and Gdependee is represented by an additional concept expression DepGdepender to handle the other case
when at least one actor is an external actors (i.e., does not
belong to the system-to-be). The dependency relation might
cause both orchestration and choreography inconsistencies,
depending on whether at least one actor is an external actor
or not. This depends on the mapping, i.e., whether an actor
is mapped to a pool (process) or to a lane (part of a process).
Thus, we will distinguish these two cases in the inconsistency detection later on. All other intentional relations can
only cause inconsistencies in the process orchestration.
RelApproveOrder≡∃ requires:ApproveOrder
⊓∃requires:CustomerProfileRetrieve ð1Þ
RelElectronicPayment≡∃requires:ElectronicPayment
⊔∃requires:InPersonPayment ð2Þ
Axiom 1 describes an AND-relation of the sibling intentional elements Approve Order and Customer Profile
Retrieve. The intentional relation of Customer Profile
Retrieve is defined equally. The AND-relation is expressed by
the concept intersection (⊓) of the concept expressions
∃ðrequires:ApproveOrder and ∃requires:CustomerProfileRetrieve.
An inclusive OR-relation between Electronic Payment or
In Person Payment is exemplified in Axiom 2, in which the
fulfillment relationship of both tasks is reflected, while both
tasks are inclusive siblings within an OR decomposition. The
relation of intention In Person Payment is defined equally.
5.2.2. Workflow relations
We represent business process models in terms of workflow (control flow) relations between activities (cf. Section
2). Algorithm 2 describes how the corresponding knowledge base ΣPM is built.
Algorithm 2. Workflow patterns ΣPM.
1: Input: Mat. process model PM ¼ðV; E; F; EAÞ
2: for all A∈A ðADVÞ do
3: OrchA≡⊤
4: end for
5: for all E∈EA then
6: if ðA1; A2Þ¼ E then
7: OrchA1 ≔OrchA1
⊓∃requires:A2
8: OrchA2 ≔OrchA2
⊓∃requires:A1
9: end if
10: end for
11: for all F∈F do
12: if
F ¼ðAND; AND; BÞ∨F ¼ðAND; IOR;BÞ∨F ¼ðAND; XOR;BÞ∨F ¼ðAND;DISC; BÞ
then
13: OrchAi ≔OrchAi
⊓⊓Bj∈B∃requires:ð⊓Ak∈BJ
Ak
Þ
14: end if
15: if F ¼ðIOR; IOR; BÞ∨F ¼ðIOR;DISC; BÞ∨F ¼ðIOR; XOR; BÞ then
16: OrchAi ≔OrchAi
⊓⊔Bj∈B∃requires:ð⊓Ak∈BJ
Ak
Þ
17: end if
18: if F ¼ðXOR; XOR; BÞ then
19: OrchAi ≔OrchAi ⊗Bj∈B∃requires:ð⊓Ak∈BJ
Ak
Þ
20: end if
21: end for
There might be an overlapping of activity relations. For
instance, the activity Money Order in Fig. 1(b) is part of an
exclusive branching fragment (internal fragment in Fig. 1
(b)) and also within a parallel branching fragment, i.e., the
activity Money Order is conjunctively and exclusively related
to other activities. Accordingly, we build orchestration
relations of activities (OrchA) as a conjunction (intersection
in DL) of activities from the different control flow patterns.
Initially, each orchestration relation concept OrchA is
defined as equivalent to the universal concept ⊤ (line 3).
In lines 5–10, the sequential control flow relations (in
both directions) of an activity A are covered by restricting
the concept OrchA. Since gateways do not realize intentions,
they are transparent in the representation (cf. Definition 3).
Afterwards, relations within fragments F are considered,
whereby only those fragments that start and end with a
gateway are relevant. Each branch B∈B of a fragment F∈F is
a set of atomic activities. In lines 12–14, the algorithm
restricts the concept definitions OrchAi
, in which Ai are those
activities that appear in parallel branches. We use an
intersection between activity sets of sibling branches, indicating the conjunctive relationship between sibling activities in parallel branches. From a logical point of view, we
treat different closing gateways (multiple merge, synchronization and discriminator) equally. Branching relations do
not impose restrictions on activities within the same branch
in a fragment (in contrast to the sequence pattern). Thus,
we describe all activities of the same branch by a concept
intersection, independent of the kind of branching.
Multi-choices are treated in lines 15–17, including synchronizing merge, simple merge and discriminator. Logically, activities of sibling branches are connected by a
concept union. In case of exclusive branching (lines 18–
20), the concept definitions OrchAi
contain a further restriction that allows only the execution of one branch.
Axiom 3 depicts a part of the control flow relation of
activity Credit Card Checking (cf. Fig. 1(b)). The activity is
part of a choice fragment, i.e., either activity Credit Card
Checking or Customer Records Checking can be executed, or
even both. This relation is represented by a concept union in
the first line of the axiom. The second line of the axiom
covers sequential relations of the activity Credit Card
Checking to its successor Approve Order.
Axiom 4 describes the relation of activity Fraud Detection, as member of a parallel branch, in which activity Apply
Discount and the internal fragment with activities Credit
Card, Debit Card and Money Order. The relation concept also
covers predecessor (Verify Payment Information) and
Table 2
Constructs and notations in DL and FOL syntax.
Construct DL syntax FOL syntax
Atomic concept CCðxÞ
Atomic role RRðx; yÞ
Concept incl. C⊑D ∀x:CðxÞ-DðxÞ
Role inclusion R⊑S ∀x; y:
Rðx; yÞ-Sðx; yÞ
Concept union C1⊔⋯⊔Cn C1ðxÞ∨⋯∨CnðxÞ
Concept int. C1⊓⋯⊓Cn C1ðxÞ∧⋯∧CnðxÞ
Concept neg. :C :CðxÞ
Universal quant. ∀P:C ∀y:ðPðx; yÞ-CðyÞÞ
Exist. quant. ∃P:C ∃y:ðPðx; yÞ∧CðyÞÞ
92 G. Gröner et al. / Information Systems 43 (2014) 83–99
successor (Build And Package Order) activities.
OrchCreditCardChecking≡ð∃ requires:CreditCardChecking
⊔∃requires:CustomerRecordsCheckingÞ
⊓∃requires:ApproveOrder ð3Þ
OrchFraudDetection≡ð∃ requires:FraudDetection
⊓∃ requires:ApplyDiscount
⊓ð∃requires:CreditCard
⊗∃requires:DebitCard
⊗∃requires:MoneyOrderÞÞ
∃ requires:VerifyPaymentInformation
⊓∃requires:BuildAndPackageOrder ð4Þ
5.2.3. Process choreography
The workflow relations of a single business process
model PM are already represented in the knowledge base
ΣPM. Finally, a representation of the process choreography
between a set of processes is needed.
Algorithm 3 depicts the axioms of a process choreography Ch in the knowledge base ΣCh.
Algorithm 3. Choreography knowledge base ΣCh.
1: Input: Process Choreography Ch ¼ðP;MÞ
2: for all Proc∈P do
3: for all A∈activitiesðProcÞ do
4: ChorA≡⊤
5: end for
6: end for
7: for all ðA1; A2Þ∈M do
8: ChorA1 ≔ChorA1
⊓∃requires:A2
9: ChorA2 ≔ChorA2
⊓∃requires:A1
10: end for
For each activity, a concept ChorA is initialized (line 4). If
there is a message exchange between activities A1 and A2,
expressed by ðA1; A2Þ∈M, we extend the concept definitions
ChorA1
and ChorA2
by referring to the other activity. According to Section 2, a message exchange is only allowed
between activities of different processes. In the concept
definitions of the choreography relations, the same DL role
requires is used as for the workflow relations that describe
the process orchestration.
5.2.4. Mapping between intentions and activities
Besides intentional relations and workflow patterns, we
have to represent the realization of tasks by the corresponding activities in terms of mappings in the knowledge base
ΣM. A mapping is described as a concept equivalence in the
knowledge base. If there is a mapping mðG; AÞ from a task G
to an activity A, we represent the mapping by an axiom
A≡G. In both models, we use the same role requires in order
to allow for a comparison of relations of both models.
6. Detection of realization inconsistencies
Mappings describe the implementation of a task (intentional element) by an activity. As a consequence, it is
expected that the relations of a task and its corresponding
activity are not contradicting. The detection of orchestration
and choreography realization inconsistencies can be
checked independently from each other, by comparing the
intentional relations of the goal model with the process
orchestration of a particular business process model and
with the choreography of a set of interacting business
processes.
6.1. Process orchestration inconsistencies
For a given mapping mðG; AÞ, we compare the corresponding workflow patterns WF of activity A and the intentional
relations IR of intentional element G in order to test whether
these relations might cause any inconsistencies. We distinguish between the two introduced kinds of inconsistencies
and no inconsistency, which is also called realization equivalence, according to the overview of inconsistencies of Table 1
(Section 4). Relations of G and A are represented by concepts
RelG and OrchA in the knowledge base.
In the knowledge base, this is reflected by the combination
of the concepts RelG and OrchA as follows: (i) A strong
inconsistency means that there cannot be any execution
combination of activities that fulfills the intentional relations
of the corresponding intentional elements. In the DL sense,
the intersection of both concepts RelG⊓OrchA is unsatisfiable,
i.e., the intersection RelG⊓OrchA cannot have a common
individual. (ii) A potential inconsistency indicates that there
might be an execution of activities, in which the corresponding intentional relation is not fulfilled. In this case, the
intersection RelG⊓OrchA is satisfiable, i.e., there can exist
common individuals of both concepts. (iii) The intentional
relations of G and the workflow patterns of activity A are
called realization equivalent if all execution combinations that
involve activity A lead to a fulfillment of the intentional
relations of G. This is true if the subsumption OrchA⊑RelG
holds. Logically, this means that OrchA implies RelG.
In order to check these three different cases, we introduce the following validation concepts. The concept Valid✓
is defined as :OrchA⊔RelG to encode the subsumption test
OrchA⊑RelG. Thus, OrchA is subsumed by RelG if Valid✓
≡:
OrchA⊔RelG is equivalent to the universal concept ⊤. This
indicates the realization equivalence. A concept Valid7 is
defined as the intersection OrchA⊓RelG to test whether there
is a potential or a strong inconsistency, i.e., if Valid7 is
satisfiable there is a potential inconsistency, otherwise a
strong inconsistency. Formally, the knowledge base ΣO is
obtained as described in Definition 10.
Definition 10 (Orchestration knowledge base ΣO). The
knowledge base ΣO≔ΣGM∪ΣPM∪ΣM is extended as follows:
For each mapping mðG; AÞ from an intentional element G to
an activity A, which is represented in ΣM by an axiom G≡A,
we insert the following axioms into ΣO:
 Valid✓
A;G≡:OrchA⊔RelG
 Valid7
A;G≡OrchA⊓RelG
Given the orchestration knowledge base, we get the
validation result en passant. Classifying the validation
G. Gröner et al. / Information Systems 43 (2014) 83–99 93
concepts Valid✓
A;G
and Valid7
A;G
of the knowledge base ΣO
leads to the following results:
1. If Valid✓
A;G
is classified equal to the universal concept ⊤,
we can guarantee the realization equivalence of the
workflow patterns over activity A and the intentional
relations over intentional element (task) G.
2. In the other case, there is either a strong inconsistency or
a potential inconsistency. This is indicated by the classification of the concept Valid7
A;G
. If Valid7
A;G
is classified as a
subconcept of the bottom concept ⊥ (i.e., Valid7
A;G⊑⊥),
there is a strong inconsistency, otherwise (i.e.,
Valid7
A;G⋢⊥) there is a potential inconsistency.
6.2. Process choreography inconsistencies
In case a goal model GM is mapped to multiple processes
and there are dependencies between actors in the goal
model, we have to check whether there is a choreography
inconsistency, according to Definition 9. Like for the orchestration inconsistencies, the mappings between intentions
and activities are represented in the mapping knowledge
base ΣM. Intentional relations are covered in the goal model
knowledge base ΣGM. Additionally, choreography between
processes of mapped activities need to be considered. We
obtain our knowledge base ΣC as described in Definition 11.
Definition 11 (Choreography knowledge base ΣC). The
knowledge base ΣC≔ΣGM∪ΣCh∪ΣM is extended as follows:
For each mapping mðG; AÞ from an intentional element G to
an activity A, which is represented in ΣM by an axiom G≡A,
we insert the following axiom into ΣC:
 ValidChor
G;A ≡:RelG⊔ChorA
From a logical point of view, this is also an implication as
for the potential inconsistency detection in Section 6.1. The
intentional relations RelG imply the choreography relation
ChorA. We expect that a dependency between intentions (in
the scope of different external actors) is also covered in the
choreography ChorA of the corresponding mapped activities.
The implication is not satisfied, which is tested in the
validation later on, if there is a dependency relation but
not the corresponding message exchange (choreography).
Otherwise, the implication holds. In the DL representation,
this means that the concept ValidChor
G;A
is equal to the
universal concept ⊤. Therefore, we detect a choreography
inconsistency if a concept ValidChor
G;A
is not equal to the
universal concept ⊤ after a concept classification by a
reasoner.
7. Proof-of concept and discussion
The first part of this section conducts a performance
evaluation of the inconsistency detection algorithms,
demonstrating the tractability of the approach for models
of realistic size. Discussions on qualitative analysis and on
the general methodology and applicability are presented
afterwards.
7.1. Performance evaluation
We conduct an evaluation by providing a proof-of concept in which workflow patterns from BPMN processes,
message exchange between processes and intentional relations from GRL goal models are represented in an OWL DL
(OWL2 DL)9
knowledge base. We analyze the reasoning
performance with different sizes of goal and process models
and with varying numbers of inconsistencies in order to test
how the validation approach is applicable for real sized
process models.
7.1.1. Experimental setting and data set
The goal of our evaluation setting is to systematically
observe how the presented validation algorithms scale with
the increase of the models and the increase of orchestration
and choreography inconsistencies. To achieve this, we apply
the simulation modeling technique by following guidelines
similar to those proposed in [34]. We selected the simulation technique, as it is commonly used in the context of
model validation and verification [35,36].
We build the models directly in the Description Logics
representation as follows:
1. We randomly generated goal models with three different sizes: (I) 200, (II) 400 and (III) 600 intentional
elements. More than 50% are tasks, i.e., intentional
elements that can be mapped to activities. The ratio/
distribution of AND, IOR and XOR decompositions is
nearly equally. Each goal model consists of 4 actors
(2 internal and 2 external actors).
2. For each goal model, 2 corresponding process models
are derived, each distributed over two lanes. For each
task, a corresponding activity is obtained, ending up
with 2 interacting process models with a total number
of (I) 100, (II) 200 and (III) 300 activities (of both
models).
 AND, IOR and XOR decompositions are “transformed” to AND, IOR and XOR workflow (control
flow) patterns.
 Positive contributions and dependencies between
internal actors lead to sequential patterns.
Due to this generation principle, the “correct” mappings between tasks and activities are already given
by the experimental process model design, instead
of manually establishing them. By this design,
we already know which activity implements
which task.
3. Finally, we modify mappings between tasks and activities such that different orchestration and choreography
inconsistencies occur. The ratio of inconsistencies varies
(for both kinds of inconsistencies) between 10%, 30%
and 50% (out of all possible inconsistencies). This
principle is based on the mutation testing technique,
9
The Web Ontology Language (OWL): http://www.w3.org/TR/owl2-o
verview/.
94 G. Gröner et al. / Information Systems 43 (2014) 83–99
originally proposed in [37] as a common fault-based
technique where simple faults are predicted.
Due to the artificial changing of mappings (third step),
we inject orchestration and choreography inconsistencies
into (originally correct) mappings between goal and process
models. Thus, we exactly know the number of possible and
existing inconsistencies in each concrete case, which paves
the way for a meaningful analysis of the performance of our
validation algorithm by varying the number of orchestration
and choreography inconsistencies.
As already mentioned, the generated process models
consist of 100, 200 and 300 activities on average. This size of
process models has been observed substantially through
existing (reference) process models in the literature [38]
such as the SAP reference process model.
The distributions of inconsistencies were chosen not
only to preserve diversity in the number of orchestration
and choreography inconsistencies, but also to have realistic
view of the number of possible inconsistencies that might
happen in real case-studies. We assume that it is rare to
have a model with 80% inconsistencies.
Table 3 shows the different settings we apply. There are
three different sizes (I)–(III) and for each size there are 9
different distributions of inconsistencies (10%, 30%, 50% of
orchestration and choreography inconsistencies), ending up
with 27 different settings. For each setting, 50 samples were
generated.
7.1.2. Evaluation results
After reasoning on the knowledge bases ΣO and ΣC, our
tool produces a list of realization equivalent mappings
(Valid✓
≡⊤) and a list of strong inconsistencies Valid7 ⊑⊥,
the remaining are known as potential inconsistencies (i.e.,
Valid7 ⋢⊥). Likewise, a list of choreography inconsistencies
(ValidChor⋢⊥) is obtained. The ontology creation is implemented with the OWL-API.10 For reasoning, we used the
Pellet reasoner.11 Our test system is a Notebook with an
Intel Core 2 Duo 8700 CPU (2.5 GHz, 4 MB cache and 2 GB
DDR2 RAM). The DL expressiveness of the knowledge bases
ΣO and ΣC is ALC.
The overall result for the different settings according to
Table 3 is depicted in Fig. 3. For each of these 27 different
settings, three different model sizes, three different numbers of orchestration inconsistencies and also three different numbers of choreography inconsistencies are applied.
The time for inconsistency detection, i.e., the time the
reasoner needs to classify the validation concepts, is the
average time for each setting (each setting consists of 50
generated models).
The observations in this experiment show that the
execution time is reasonable for realistic-size models and
for a varying number of inconsistencies. As shown in Fig. 3,
the size of the models affect the execution time. The mean
values of the execution time are: 1493.04 ms (100 activities), 2807.74 ms (200 activities) and 4394.27 (300 activities). This is a rather expected result since in the increase in
the model size also increase proportionally the number of
mappings and therefore also the number of validation
concepts (three for each mapping) that have to be classified.
We can observe that the execution time of our algorithm for
the largest experimented model with 300 activities is less
than 5.3 s.
Regarding the number of inconsistencies, the experiment demonstrates that the increase of orchestration inconsistencies increases the execution time of the algorithm.
Independent of the model size, an increase of the percentage of orchestration inconsistencies increases the validation time (with a fixed model size and a fixed number of
choreography inconsistencies) on an average by 7–9%. For
instance, from setting 10–30 to setting 10–50 the mean time
(M) increases from 1486 to 1613 ms. There is no difference
between strong and potential inconsistencies.
The increase of choreography inconsistencies has less
influence on the execution time, only a slight increase is
shown. A possible explanation is the lower number of
relationships that are covered by the concept expressions
ðChorA and ValidChorÞ since in contrast to orchestrations
(strong and potential inconsistencies) only dependencies
are covered by these concept expressions and no other
intentional relations. The increase is on an average by 3–5%.
7.1.3. Threats to validity
An experimental evaluation is always subject to different
threats that can affect the validity of the results. In the
following, we investigate possible threats to the validity of
our results. Three main factors influence the validity of our
experiments: (i) the size of the models, (ii) the distribution
of intentional (and therefore also control flow) relations and
(iii) the modeling approach. With respect to the size of the
process models, our investigation over existing publications
in process management confirmed that the sizes of generated models are aligned with the size of existing models in
the literature. The same holds for the assumed distribution
of intentional relations (AND, IOR and XOR), which is
assumed as approximately equal. Regarding the proposed
modeling and formalization approach of process models
and goal models, we may have slightly different execution
times for different formal representations. Thus, our results
cannot be generalized to other formalisms, which could be
used as alternative for implementation of the proposed
validation approach. However, the experimental results
show that employing Description Logic is a tractable means
to model and validate goal-oriented process models.
7.2. Validation exemplified
We demonstrate the validation for an excerpt of Fig. 1.
Consider the strong inconsistency for the activities Send
Receipt and Shipment.
12 They are siblings in exclusive branches
of the same fragment. In ΣPM, there are the definitions of this
10 OWL-API site: http://owlapi.sourceforge.net/.
11 Pellet reasoner site: http://clarkparsia.com/pellet/.
12 Note that we represent here the exclusive relation in the standard
Description Logics notation to ease the understanding of the inconsistency
that is caused by the existence of a positive and its negated statement
(after mapping Send Receipt to Bill Preparation (BL)).
G. Gröner et al. / Information Systems 43 (2014) 83–99 95
relation for both activities. An excerpt of the relation definition
of activity Send Receipt is shown in Axiom 5.
OrchSendReceipt≡ð∃requires:SendReceipt
⊔∃requires:ShipmentÞ
⊓:ð∃requires:SendReceipt
⊓∃requires:ShippingÞð5Þ
RelBL≡∃requires:ShipOrder
⊓∃requires:BL ð6Þ
Since Send Receipt is mapped to the intentional element
(task) Bill Preparation (BL), Send Receipt is defined as
equivalent to the concept Bill Preparation (BL) and Shipment
is equivalent to Ship Order (mapping knowledge base ΣM).
From the goal model, there is a relation definition of Bill in
ΣGM as depicted in Axiom 6.
The validation concept Valid✓
≡:OrchSendReceipt⊔RelBL is
classified by the reasoner as different from the universal
concept ⊤, i.e., we know that the realization equivalence
does not hold between Send Receipt and Bill Preparation.
The other validation concept Valid7 ≡OrchSendReceipt⊓RelBL is
classified equivalent to the bottom concept ⊥, i.e., indicating
a strong inconsistency. The same holds for Shipment.
7.3. Lessons learned
The course of our research, as presented in the previous
part of this paper, has raised the following issues for further
discussions.
Inconsistency detection and validation: The purpose of the
validation is to detect inconsistencies between mapped
elements with respect to their relationships, i.e., intentional
relations and relations of activities. This is based on a
comparison how these elements are logically related to
each other, e.g., whether their relations coincide or contradict to each other. However, an analysis whether or to
which degree the execution of an activity satisfies a particular goal from a qualitative perspective leads to further
interesting research questions, e.g., like the consideration of
soft goal satisfaction in reasoning algorithms.
Resolution of detected inconsistencies: The motivation of
the presented work is to ensure correct mappings between
two different kinds of models. We do not restrict how a
detected inconsistency can be resolved. In general, there are
three non-disjoint possible directions: (i) changing the
mappings, (ii) redesigning the process model or (iii) redesigning the goal model.
However, another application context could be the dedicated redesign of process models, in which inconsistent process
models (with respect to goal models) are redesigned until they
meet the requirements that are imposed by the goal models.
From requirements to orchestration. As mentioned in
Section 4, we restrict mappings to tasks in the goal model
since they operationalize stakeholders' goals, and therefore,
they can be directly realized by activities in a process. The
proposed modeling framework is based on these assumptions. Discussions whether hard and especially soft goals
could be mapped to activities and whether this is meaningful are outside of the scope of this paper.
From requirements to choreography: As discussed in
Section 4, the inconsistency detection does not check the
directions of dependencies and interaction patterns, since
the dependency relations in the goal model do not provide
enough information to decide the type of service interaction
patterns between the corresponding processes, we cannot
make a conclusion about the types and directions of interaction patterns based on dependency relations. Hence, we
did not make a claim in this regard. However, it is an
Table 3
Characteristics of the generated models.
No. Goal model [intentions] Process model [activities] Realizations [mappings] Choreography Inc. (%) Orchestration Inc. (%)
(I) 200 100 100 [10, 30, 50] [10, 30, 50]
(II) 400 200 200 [10, 30, 50] [10, 30, 50]
(III) 600 300 200 [10, 30, 50] [10, 30, 50]
Fig. 3. The results of the experiment.
96 G. Gröner et al. / Information Systems 43 (2014) 83–99
interesting point for further research investigations and
case studies to investigate whether and how directions of
dependencies in goal models influence the directions of
message exchange between process models.
Overall application context: Another aspect, to be discussed, is the usefulness and applicability in real applications. The key benefit of the modeling setting with goal and
process models in combination with mappings offers two
different “views” on a system/application (requirements
and process models), while ensuring the correct alignment
between modeling constructs and their logics in both views.
This allows a (partly) independent development and maintenance of a system from two viewpoints. Alternative
approaches consider the derivation of a process model from
an existing requirement model (cf. Section 8).
7.4. Rationale for Description Logics
Description Logics is an expressive language for knowledge representation. The semantics of Description Logics
provides an unambiguous interpretation of expressions in
the knowledge base that is a prerequisite for automated
reasoning. The contribution of Description Logics reasoning
to the validation is threefold:
 We use reasoning to recognize orchestration and choreography inconsistencies.
 Due to the classification of concepts and the subsumption checking between concepts, the reasoner directly
pinpoints the source of an inconsistency, i.e., which
element (intention and activity) is part of an inconsistency. This is crucial in large models, where various
mappings exist.
 We use classifications of concept expressions by the
reasoner in order to check whether a concept expression
is always satisfied, is satisfiable or is unsatisfiable. This is
exploited to distinguish between strong and potential
inconsistencies.
The strengths of Description Logics are quite efficient
reasoning procedure in practical settings and there is a well
established infrastructure for modeling and reasoning tool
support. Description Logics is a decidable subset of firstorder logic, and the theoretical exponential reasoning complexity is tractable in practical settings. In contrast to
propositional logic, Description Logics is more expressive
and allows the integration of further background knowledge like annotations of elements.
Finally, there are some weaknesses of Description Logics
and the proposed modeling approach. Certain aspects of the
models cannot be represented in standard Description Logic
models, e.g., the influence of qualitative aspects, as captured
by soft goals in which there is no clear-cut criteria to
measure the achievement, need further modeling
approaches.
Our representation relies on a structural representation of
processes and how activities are related to each other.
Temporal relations are covered by predecessor and successor
relations of activities, but without (rather powerful) temporal
(logic) representation formalisms. Due to transformations
and dedicated modeling decisions, like the introduced materialized process model representation (cf. Section 2), we can
disregard temporal constructs. However, such temporal
modeling constructs are widely used in process modeling
and the incorporation of them is another interesting aspect.
The standard Description Logics language does not support
temporal constructs. However, there are some initial
approaches that address such issues in Description Logics
[39].
8. Related work
The first group of related work considers relations
between goal models and business process models, but
not validation as it is done in our work. Transforming goal
models into business process model has been one of the
major research areas in business process management.
Lapouchnian et al. [40] use goal models for the configuration of business processes. Goal models are annotated with
control flow information and afterwards transformed into
BPEL processes. Similarly, Decreus and Poels [41] annotate
goal-oriented models in the so-called B-SCP framework
with control flow information and transform them into
BPMN skeletons. Furthermore, Frankova et al. [42] transform SI*/Secure Tropos models into skeletons of process
models in BPMN, from which they generate executable
processes in BPEL. On the other hand, Santos et al. [5]
derive goal models from existing process models. Such goal
models are then used to control variability and configuration of processes. Finally, Koliadis et al. [3] annotate activities of process models with effects, whereby these effects
serve for a comparison of a process with user goals, i.e.,
effects of activities are compared with goal fulfillment. The
basic principle is to reflect changes from an i
⋆ model to a
BPMN model and vice versa.
The second group of related research is more related to
our work, where some kind of verification between requirements and business process models is investigated. In this
line of related research, Kazhamiakin et al. [29] propose a
methodology that provides a set of high-level mapping
rules to produce process models in BPEL from Tropos
models. In order to enable requirements driven verification
of process models, they employ the formal Tropos language
in combination with temporal constraints. While their work
and their tools support several types of formal analyses for
consistency checks in the requirements models as well as in
the process model, they do not focus on validation of the
inconsistencies across models, i.e., those cases where both
models are consistent but the mapping align elements with
contradicting relationships, as it is done in our work.
Soffer and Wand [43,44] propose a generic theory-based
process modeling framework based on Bunge's ontology for
a goal-driven analysis of process models to check goal
reachability in process models. They define a goal as stable
state, which indicates the termination of the process. Their
framework defines a set of assumptions and parameters
(based on goal relations and workflow relations), which are
used to check if a set of workflow patterns ensure that
processes can always reach their goals. Hence, using these
parameters, it is possible to identify sets of valid and invalid
design decision with respect to business goals. They also
G. Gröner et al. / Information Systems 43 (2014) 83–99 97
suggest appropriate redesign actions to eliminate invalid
designs. Our approach follows the similar objective, but we
use DL based reasoning to identify inconsistencies
automatically.
Pistore et al. [45] introduce an approach to design and
verify web services. The requirements are also specified
using goal models or more precisely Formal Tropos models.
BPEL4WS is used for business processes. Their approach
enables verification of whether linear time temporal constraints specified in goal models hold in business process
models. Mahfouz et al. [46] complement the approach of
Pistore et al. and verify if linear time first order temporal
constraints, specified in goal models, hold in message
exchange (choreography). Our approach verifies that intentional relationships in goal models are consistent with
execution and message exchange in process models. Similar
to our approach, Markovic et al. [47] introduce an ontology
based specification of goal models and their relationships to
process models. However, potential inconsistencies and
validation of goals in process models are not discussed.
Liaskos et al. [48] also introduce an approach for goal
based customization of workflows. A family of processes is
extended with the notation for partial temporal ordering of
goals. A particular member of the family is specified by
constraints with the means of linear temporal logic operators. However, this approach does not take into consideration business process logic embedded in the realization of
business processes, which is the focus of our validation. La
Rosa et al. [49] propose a questionnaire based customization of configurable business processes represented in CYAWL [50]. They use a Petri-net based reasoner [51] to
preserve the correctness during process configuration.
Variability patterns of La Rosa et al.'s work are narrower
than patterns introduced in this work.
9. Conclusion
In this paper, we have presented a novel approach for
handling inconsistencies in mappings between goal models
and business process models. We distinguish between two
kinds of inconsistencies. (i) Inconsistencies of the process
orchestration are caused by contradicting control flow
relationships between activities in the business process
model and between relationships of user intentions in the
goal model. (ii) Choreography inconsistencies occur if
dependencies among actors in the goal model are not
reflected by message exchange between the corresponding
activities in the business process models. By automatically
identifying these inconsistencies, we are able to detect
executable processes that did not meet user requirements
or lead to undesired executions. Additionally, we allow for
goal models and business process models to evolve
independently.
Our contribution extends the body of knowledge in the
field by considering these mapping as first-class citizens
along with goal models and business process models. We
plan to extend this approach in combination with our
existing work on configuration of business process families
[52] to provide a complete solution for software product
lines that use goal models, feature models and process
models as main artifacts.
References
[1] M. Dumas, J. Recker, M. Weske, Management and engineering of
process-aware information systems: introduction to the special issue,
Information Systems 37 (2) (2012) 77–79.
[2] A. Pourshahid, D. Amyot, L. Peyton, S. Ghanavati, P. Chen, M. Weiss, A.J.
Forster, Business process management with the user requirements
notation, Electronic Commerce Research, 9 (4) (2009) 269–316.
[3] G. Koliadis, A. Vranesevic, M. Bhuiyan, A. Krishna, A.K. Ghose,
Combining i* and BPMN for business process model lifecycle management, in: Business Process Management Workshops, Lecture Notes in
Computer Science, vol. 4103, Springer, 2006, pp. 416–427.
[4] P. Kueng, P. Kawalek, Goal-based business process models: creation
and evaluation, Business Process Management Journal (1997) 17–38.
[5] E. Santos, J. Castro, J. Sánchez, O. Pastor, A goal-oriented approach for
variability in BPMN, in: Workshop em Engenharia de Requisitos, WER,
2010.
[6] B. Andersson, M. Bergholtz, A. Edirisuriya, T. Ilayperuma, P. Johannesson, A declarative foundation of process models, in: Advanced
Information Systems Engineering, Springer, 2005, pp. 233–247.
[7] M. Bergholtz, P. Jayaweera, P. Johannesson, P. Wohed, A pattern and
dependency based approach to the design of process models, in:
Conceptual Modeling—ER, 23rd International Conference, Lecture
Notes in Computer Science, vol. 3288, Springer, 2004, pp. 724–739.
[8] N. Russel, A. ter Hofstede, W. van der Aalst, N. Mulyar, Workflow
Control-Flow Patterns: A Revised View, Technical Report, BPM Center
Report BPM-06-22, BPMcenter.org, 2006.
[9] E. Yu, P. Giorgin, N. Maiden, J. Mylopoulos (Eds.), Social Modeling for
Requirements Engineering, MIT, 2011.
[10] P. Giorgini, J. Mylopoulos, R. Sebastiani, Goal-oriented requirements
analysis and reasoning in the TROPOS methodology, Engineering
Applications of Artificial Intelligence 18 (2005) 159–171.
[11] L. Chung, B.A. Nixon, E. Yu, J. Mylopoulos, Non-Functional Requirements in Software Engineering, first ed. Springer, 1999.
[12] A. van Lamsweerde, Requirements Engineering: From System Goals to
UML Models to Software Specifications, Wiley, 2009.
[13] ITU-T, Recommendation Z.151 (09/08): User Requirements Notation
(URN) Language Definition, 2008.
[14] D. Amyot, S. Ghanavati, J. Horkoff, G. Mussbacher, L. Peyton, E. Yu,
Evaluating goal models within the goal-oriented requirement language,
International Journal on Intelligent Systems 25 (2010) 841–877.
[15] E.S.K. Yu, J. Mylopoulos, Understanding “Why” in software process
modelling, analysis and design, in: Proceedings of the 16th International Conference on Software Engineering, ICSE '94, IEEE Computer
Society Press, Los Alamitos, CA, USA, 1994, pp. 159–168. URL: 〈http://
dl.acm.org/citation.cfm?id=257734.257757〉.
[16] C. Ouyang, M. Dumas, W.M.P.V.D. Aalst, A.H.M.T. Hofstede, J. Mendling,
From business process models to process-oriented software systems,
ACM Transactions on Software Engineering and Methodology 19
(2009) 2:1–2:37.
[17] C. Peltz, Web services orchestration and choreography, Computer 36
(10) (2003) 46–52.
[18] P. Feiler, W. Humphrey, Software process development and enactment: concepts and definitions, in: 2nd International Conference on
the Software Process, Continuous Software Process Improvement,
1993, pp. 28 –40. http://dx.doi.org/10.1109/SPCON.1993.236824.
[19] W. van der Aalst, A. ter Hofstede, B. Kiepuszewski, A. Barros, Workflow
patterns, in: Distributed and Parallel Databases, 2003.
[20] W.M.P. van der Aalst, A.H.M. ter Hofstede, YAWL: yet another workflow language, Information Systems 30 (4) (2005) 245–275.
[21] O.M.G. (OMG), Business Process Model and Notation (BPMN), Version
2.0, Technical Report formal/2011-01-03, 2011. URL: 〈http://www.
omg.org/spec/BPMN/2.0/PDF/〉.
[22] B. Kiepuszewski, A.H.M.t. Hofstede, C. Bussler, On structured workflow modelling, in: 12th International Conference on Advanced
Information Systems Engineering, CAiSE, Springer, London, UK,
2000, pp. 431–445.
[23] J.M. Küster, C. Gerth, A. Förster, G. Engels, Detecting and resolving
process model differences in the absence of a change log, in: Business
Process Management, 6th International Conference, BPM, Lecture
Notes in Computer Science, vol. 5240, Springer, 2008, pp. 244–260.
[24] J. Vanhatalo, H. Völzer, F. Leymann, Faster and more focused controlflow analysis for business process models through SESE decomposition, in: Service-Oriented Computing—ICSOC, Lecture Notes in Computer Science, vol. 4749, Springer, 2007, pp. 43–55.
[25] R. Johnson, D. Pearson, K. Pingali, The program structure tree:
computing control regions in linear time, in: PLDI, 1994, pp. 171–185.
[26] S. White, D. Miers, BPMN Modeling and Reference Guide, Future
Strategies Inc., 2008.
98 G. Gröner et al. / Information Systems 43 (2014) 83–99
[27] S. Liaskos, A. Lapouchnian, Y. Yu, E. Yu, J. Mylopoulos, On goal-based
variability acquisition and analysis, in: Requirements Engineering,
14th IEEE International Conference, 2006, pp. 79–88.
[28] A. Tarski, The Semantic Conception of Truth and the Foundations of
Semantics, Philosophy and Phenomenological Research, 1944.
[29] R. Kazhamiakin, M. Pistore, M. Roveri, A framework for integrating
business processes and business requirements, in: Proceedings of
IEEE EDOC Conference, 2004, pp. 9–20.
[30] A. Mahfouz, L. Barroca, R.C. Laney, B. Nuseibeh, From organizational
requirements to service choreography, in: World Conference on
Services-I, IEEE, 2009, pp. 546–553.
[31] A. Barros, M. Dumas, A.T. Hofstede, Service Interaction Patterns:
Towards a Reference Framework for Service-based Business Process
Interconnection, Technical Report, 2005.
[32] G. Decker, F. Puhlmann, Extending BPMN for modeling complex
choreographies, in: Proceedings of the 2007 OTM Confederated
International Conference on On the Move to Meaningful Internet
Systems: CoopIS, DOA, ODBASE, GADA, and IS, vol. Part I, OTM'07,
Springer-Verlag, Berlin, Heidelberg, 2007, pp. 24–40. URL: 〈http://dl.
acm.org.proxy.lib.sfu.ca/citation.cfm?id=1784607.1784614〉.
[33] F. Baader, D. Calvanese, D.L. McGuinness, D. Nardi, P.F. Patel-Schneider
(Eds.), The Description Logic Handbook, second ed. Cambridge University Press, 2007.
[34] M.I. Kellner, R.J. Madachy, D.M. Raffo, Software process simulation
modeling: Why? What? Journal of Systems and Software 46 (1999)
91–105.
[35] J. Guo, Y. Wang, P. Trinidad, D. Benavides, Consistency maintenance
for evolving feature models, Expert Systems with Applications 39 (5)
(2012) 4987–4998.
[36] S. Sadiq, M. Orlowska, W. Sadiq, Specification and validation of
process constraints for flexible workflows, Information Systems 30
(5) (2005) 349–378.
[37] R.A. DeMillo, R.J. Lipton, F.G. Sayward, Hints on test data selection:
help for the practicing programmer, Computer 11 (4) (1978) 34–41,
http://dx.doi.org/10.1109/C-M.1978.218136. URL: 〈http://dx.doi.org.
proxy.lib.sfu.ca/10.1109/C-M.1978.218136〉.
[38] J. Mendling, Metrics for Process Models: Empirical Foundations of
Verification, Error Prediction, and Guidelines for Correctness, first ed.
Springer, 2008.
[39] C. Lutz, F. Wolter, M. Zakharyaschev, Temporal description logics: a
survey, in: 15th International Symposium on Temporal Representation and Reasoning (TIMETIME), IEEE Computer Society, 2008, pp.
3–14.
[40] A. Lapouchnian, Y. Yu, J. Mylopoulos, Requirements-driven design and
configuration management of business processes, in: Business Process
Management, 5th International Conference, BPM, Lecture Notes in
Computer Science, vol. 4714, Springer, 2007, pp. 246–261.
[41] K. Decreus, G. Poels, A goal-oriented requirements engineering
method for business processes, in: CAiSE Forum, Lecture Notes in
Computer Science, vol. 72, 2010, pp. 29–43.
[42] G. Frankova, G. Frankova, F. Massacci, F. Massacci, M. Seguran, M.
Seguran, From Early Requirements Analysis towards Secure Workflows, Technical Report TR DIT-07-036, University of Trento, 2007.
[43] P. Soffer, Y. Wand, Goal-driven analysis of process model validity, in:
Advanced Information Systems Engineering, Lecture Notes in Computer Science, vol. 3084, Springer, 2004, pp. 229–319. URL: 〈http://dx.
doi.org/10.1007/978-3-540-25975-6_37〉.
[44] P. Soffer, Y. Wand, M. Kaner, Semantic analysis of flow patterns in
business process modeling, in: 5th International Conference on
Business Process Management (BPM), Lecture Notes in Computer
Science, vol. 4714, Springer, 2007, pp. 400–407.
[45] M. Pistore, M. Roveri, P. Busetta, Requirements-driven verification of
web services, Electronic Notes in Theoretical Computer Science 105
(2004) 95–108.
[46] A. Mahfouz, L. Barroca, R.C. Laney, B. Nuseibeh, Requirements-driven
collaborative choreography customization, in: ICSOC/ServiceWave,
2009, pp. 144–158.
[47] I. Markovic, M. Kowalkiewicz, Linking business goals to process
models in semantic business process modeling, in: The 12th International IEEE Enterprise Distributed Object Computing Conference,
EDOC '08, 2008, pp. 332–338.
[48] S. Liaskos, M. Litoiu, M.D. Jungblut, J. Mylopoulos, Goal-based behavioral customization of information systems, in: 23rd International
Conference on Advanced Information Systems Engineering, CAiSE,
Springer, 2011, pp. 77–92.
[49] M. La Rosa, W. van der Aalst, M. Dumas, A. ter Hofstede,
Questionnaire-based variability modeling for system configuration,
Software and Systems Modeling 8 (2009) 251–274.
[50] F. Gottschalk, W.M.P. van der Aalst, M.H. Jansen-Vullers, M.L. Rosa,
Configurable workflow models, International Journal on Cooperative
Information Systems 17 (2) (2008) 177–221.
[51] W. van der Aalst, M. Dumas, F. Gottschalk, A. ter Hofstede, M. La Rosa,
J. Mendling, Preserving correctness during business process model
configuration, Formal Aspects of Computing 22 (2010) 459–482.
[52] G. Gröner, C. Wende, M. Bošković, F.S. Parreiras, T. Walter, F. Heidenreich, D. Gašević, S. Staab, Validation of families of business processes,
in: 23rd International Conference on Advanced Information Systems
Engineering (CAiSE), Lecture Notes in Computer Science, vol. 6741,
Springer, 2011, pp. 551–565.
G. Gröner et al. / Information Systems 43 (2014) 83–99 99